% -*- coding: utf-8 -*-
% !TEX program = xelatex
\documentclass[plain,noanswer,medmath]{../../template/jnuexam}
\usepackage{../../template/kysx}

\begin{document}


\examtitle{2015}{3}


\loadproblems{1}{2015P1}%载入试卷一
\loadproblems{2}{2015P2}%载入试卷二


\makepart{选择题}{1～8小题，每小题4分，共32分}


\begin{problem}
设$\{x_n\}$是数列，下列命题中不正确的是\pickout{}
\begin{abcd}
\item 若$\lim_{n\to\infty} x_n = a$，则$\lim_{n\to\infty} x_{2n} = \lim_{n \to \infty} x_{2n + 1} = a$．
\item 若$\lim_{n\to\infty} x_{2n} = \lim_{n\to\infty} x_{2n + 1} = a$，则$\lim_{n \to \infty} x_n = a$．
\item 若$\lim_{n\to\infty} x_n = a$，则$\lim_{n\to\infty} x_{3n} = \lim_{n \to \infty} x_{3n + 1} = a$．
\item 若$\lim_{n\to\infty} x_{3n} = \lim_{n\to\infty} x_{3n + 1} = a$，则$\lim_{n \to \infty} x_n = a$．
\end{abcd}
\end{problem}


\useproblem{1}{1}{1}%【同试卷一第一(1)题】


\begin{problem}
设$D = \left\{(x,y) \middle| x^2 + y^2 \le 2x, x^2 + y^2 \le 2y \right\}$，
函数$f\left(x,y \right)$在$D$上连续，\goodbreak 则$\iint_D f\left(x,y \right)\dx\dy =$\pickout{}
\begin{abcd}
\item $\int_0^{\frac{\pi}{4}} \d\theta \int_0^{2\cos \theta} f\left(r\cos \theta ,r\sin \theta \right) r\dr 
 + \int_{\frac{\pi}{4}}^{\frac{\pi}{2}}\d\theta \int_0^{2\sin\theta}f(r\cos\theta,r\sin\theta)r\dr$．
\item $\int_0^{\frac{\pi}{4}} \d\theta \int_0^{2\sin \theta} f\left(r\cos \theta ,r\sin \theta \right) r\dr 
 + \int_{\frac{\pi}{4}}^{\frac{\pi}{2}} \d\theta \int_0^{2\cos \theta} f(r\cos\theta,r\sin\theta) r\dr$．
\item $2\int_0^1 \dx \int_{1 - \sqrt{1 - x^2}}^x f(x,y)\dy$．
\item $2\int_0^1 \dx \int_x^{\sqrt{2x - x^2}} f(x,y)\dy$．
\end{abcd}
\end{problem}


\begin{problem}
下列级数中发散的是\pickout{}
\begin{abcd}
\item $\sum_{n = 1}^{\infty}\frac{n}{3^n}$．
\item $\sum_{n = 1}^{\infty}\frac{1}{\sqrt{n}}\ln \left(1 + \frac{1}{n} \right)$．
\item $\sum_{n = 2}^{\infty}\frac{(- 1)^n + 1}{\ln n}$．
\item $\sum_{n = 1}^{\infty}\frac{n!}{n^n}$．
\end{abcd}
\end{problem}


\useproblem{1}{1}{5}%【同试卷一第一(5)题】


\useproblem{1}{1}{6}%【同试卷一第一(6)题】


\useproblem{1}{1}{7}%【同试卷一第一(7)题】


\begin{problem}
设总体$X\sim B\left(m,\theta \right)$，$X_1,X_2, \ldots ,X_n$为来自该总体的简单随机样本，
$\overline{X}$为样本均值，则$E\left[\sum_{i = 1}^n \left(X_i - \overline{X}\right)^2 \right] =$\pickout{}
\begin{abcd}
\item $(m - 1)n\theta \left(1 - \theta \right)$．
\item $m(n - 1)\theta \left(1 - \theta \right)$．
\item $(m - 1)(n - 1)\theta \left(1 - \theta \right)$．
\item $mn\theta \left(1 - \theta \right)$．
\end{abcd}
\end{problem}


\makepart{填空题}{9～14小题，每小题4分，共24分}


\useproblem{1}{2}{9}%【同试卷一第二(9)题】


\useproblem{2}{2}{11}%【同试卷二第二(11)题】


\useproblem{2}{2}{13}%【同试卷二第二(13)题】


\useproblem{2}{2}{12}%【同试卷二第二(12)题】


\useproblem{2}{2}{14}%【同试卷二第二(14)题】


\useproblem{1}{2}{14}%【同试卷一第二(14)题】


\makepart{解答题}{15～23题，共94分}


\useproblem[points=10]{1}{3}{15}%【同试卷一第三(15)题】


\useproblem[points=10]{2}{3}{18}%【同试卷二第三(18)题】


\begin{problem}[points=10]
为了实现利润的最大化，厂商需要对某商品确定其定价模型，设$Q$为该商品的需求量，$P$为价格，
$MC$为边际成本，$\eta$为需求弹性$(\eta > 0)$．
\begin{enumerate}
\item 证明定价模型为$P = \frac{MC}{1 - \frac{1}{\eta}}$；
\item 若该商品的成本函数为$C(Q)=1600+Q^2$，需求函数为$Q=40-P$，试由(Ⅰ)中的定价模型确定此商品的价格．
\end{enumerate}
\end{problem}


\useproblem[points=10]{1}{3}{16}%【同试卷一第三(16)题】


\useproblem[points=10]{1}{3}{18}%【同试卷一第三(18)题】


\useproblem[points=11]{2}{3}{22}%【同试卷二第三(22)题】


\useproblem[points=11]{1}{3}{21}%【同试卷一第三(21)题】


\useproblem[points=11]{1}{3}{22}%【同试卷一第三(22)题】


\useproblem[points=11]{1}{3}{23}%【同试卷一第三(23)题】


\end{document}
